\documentclass[UTF8,a4paper,11pt]{ctexart}

\usepackage{amsmath,amssymb,bm,mathtools}
\usepackage{geometry}
\usepackage{hyperref}

\geometry{margin=2.5cm}
\hypersetup{
  colorlinks=true,
  linkcolor=blue,
  citecolor=blue,
  urlcolor=blue
}

\newcommand{\sat}{\operatorname{sat}}
\newcommand{\sgn}{\operatorname{sgn}}

\title{分数阶滑模控制双积分系统案例}
\author{}
\date{}

\begin{document}
\maketitle

\section{问题描述}

考虑带有匹配扰动的双积分系统
\begin{equation}
  \dot{x}_1(t)=x_2(t),\qquad
  \dot{x}_2(t)=u(t)+d(t),
  \label{eq:plant}
\end{equation}
其中 $x_1$ 为位置状态，$x_2$ 为速度状态，$u$ 为控制输入，$d(t)$ 为有界匹配扰动。假设扰动满足
\begin{equation}
  |d(t)|\le \bar d,\qquad \bar d>0.
\end{equation}
控制目标是在扰动有界的情况下，使状态 $x_1(t),x_2(t)$ 收敛到原点附近；若采用理想符号函数，则实现渐近收敛并具有有限时间到达滑模面的性质。

\section{分数阶算子}

令 $0<\alpha<1$。本文采用 Riemann--Liouville 分数阶积分
\begin{equation}
  {}_0 I_t^\mu f(t)
  =
  \frac{1}{\Gamma(\mu)}
  \int_0^t (t-\tau)^{\mu-1} f(\tau)\,d\tau,
  \qquad \mu>0,
\end{equation}
以及 Riemann--Liouville 分数阶导数
\begin{equation}
  {}_0^{\mathrm{RL}}D_t^\alpha f(t)
  =
  \frac{d}{dt}\,{}_0 I_t^{1-\alpha}f(t).
\end{equation}
当初值项处理为零，或将初值补偿显式并入控制律时，可用 Caputo 导数替换上述导数。本文的滑模面使用分数阶积分项，因此控制律中自然出现 ${}_0^{\mathrm{RL}}D_t^\alpha x_1$。

\section{分数阶滑模面设计}

定义分数阶滑模变量
\begin{equation}
  s(t)=x_2(t)+\lambda\,{}_0 I_t^{1-\alpha}x_1(t),
  \qquad \lambda>0.
  \label{eq:surface}
\end{equation}
该滑模面不是单纯的整数阶线性组合，而是把位置状态的历史记忆项引入到滑模变量中。对 \eqref{eq:surface} 求导可得
\begin{equation}
  \dot{s}(t)
  =
  \dot{x}_2(t)+\lambda\,{}_0^{\mathrm{RL}}D_t^\alpha x_1(t)
  =
  u(t)+d(t)+\lambda\,{}_0^{\mathrm{RL}}D_t^\alpha x_1(t).
  \label{eq:sdot-open}
\end{equation}

\section{控制律}

设计分数阶滑模控制律为
\begin{equation}
  u(t)
  =
  -\lambda\,{}_0^{\mathrm{RL}}D_t^\alpha x_1(t)
  -k s(t)
  -\rho\,\sgn(s(t)),
  \label{eq:control-sign}
\end{equation}
其中 $k>0$，$\rho>\bar d$。将 \eqref{eq:control-sign} 代入 \eqref{eq:sdot-open}，得到闭环滑模变量动力学
\begin{equation}
  \dot{s}(t)
  =
  -k s(t)-\rho\,\sgn(s(t))+d(t).
  \label{eq:sdot-closed}
\end{equation}

为降低抖振，实际实现中可使用边界层饱和函数
\begin{equation}
  \sat\!\left(\frac{s}{\phi}\right)
  =
  \begin{cases}
    \sgn(s), & |s|>\phi,\\[2mm]
    s/\phi, & |s|\le \phi,
  \end{cases}
  \qquad \phi>0,
\end{equation}
并令
\begin{equation}
  u(t)
  =
  -\lambda\,{}_0^{\mathrm{RL}}D_t^\alpha x_1(t)
  -k s(t)
  -\rho\,\sat\!\left(\frac{s(t)}{\phi}\right).
  \label{eq:control-sat}
\end{equation}
理想稳定性证明通常先基于符号函数 \eqref{eq:control-sign} 给出；饱和函数 \eqref{eq:control-sat} 则对应实际系统中的最终有界稳定。

\section{到达阶段稳定性分析}

取 Lyapunov 函数
\begin{equation}
  V_s(t)=\frac{1}{2}s^2(t).
\end{equation}
由 \eqref{eq:sdot-closed} 可得
\begin{align}
  \dot V_s
  &=
  s\dot{s} \notag\\
  &=
  -k s^2-\rho |s|+s d(t) \notag\\
  &\le
  -k s^2-\rho |s|+|s|\,|d(t)| \notag\\
  &\le
  -k s^2-(\rho-\bar d)|s|.
  \label{eq:reaching}
\end{align}
由于 $k>0$ 且 $\rho>\bar d$，有
\begin{equation}
  \dot V_s\le 0.
\end{equation}
更进一步，令 $\eta=\rho-\bar d>0$，则
\begin{equation}
  \dot V_s\le -\eta |s|
  =
  -\eta\sqrt{2V_s}.
\end{equation}
因此 $s(t)$ 可在有限时间内到达滑模面 $s=0$，到达时间满足估计
\begin{equation}
  T_r
  \le
  \frac{|s(0)|}{\rho-\bar d}.
\end{equation}
该结论说明，只要切换增益 $\rho$ 大于扰动上界，滑模面具有鲁棒可达性。

\section{滑动阶段稳定性分析}

当系统进入理想滑模面 $s(t)=0$ 后，由 \eqref{eq:surface} 得到
\begin{equation}
  x_2(t)+\lambda\,{}_0 I_t^{1-\alpha}x_1(t)=0.
\end{equation}
结合 $\dot{x}_1=x_2$，滑动阶段的等效动力学为
\begin{equation}
  \dot{x}_1(t)
  =
  -\lambda\,{}_0 I_t^{1-\alpha}x_1(t).
  \label{eq:sliding-dynamics}
\end{equation}
对 \eqref{eq:sliding-dynamics} 作 Laplace 变换。忽略初值项时，有
\begin{equation}
  pX_1(p)+\lambda p^{-(1-\alpha)}X_1(p)=0.
\end{equation}
对应的特征方程为
\begin{equation}
  p+\lambda p^{\alpha-1}=0,
\end{equation}
即
\begin{equation}
  p^{2-\alpha}+\lambda=0.
  \label{eq:char}
\end{equation}
记
\begin{equation}
  \beta=2-\alpha.
\end{equation}
由于 $0<\alpha<1$，所以 $1<\beta<2$。特征方程 \eqref{eq:char} 等价于标量分数阶系统
\begin{equation}
  {}_0^{\mathrm{C}}D_t^\beta x_1(t)
  =
  -\lambda x_1(t)
\end{equation}
的特征方程。由 Matignon 稳定性判据，阶次为 $\beta$ 的线性分数阶系统
$ {}_0^{\mathrm{C}}D_t^\beta x=Ax $ 稳定的充分必要条件为矩阵 $A$ 的所有特征值
$\mu_i$ 满足
\begin{equation}
  |\arg(\mu_i)|>\frac{\beta\pi}{2}.
\end{equation}
本例中等效系统矩阵为标量 $A=-\lambda$，其唯一特征值为
\begin{equation}
  \mu=-\lambda.
\end{equation}
因为 $\lambda>0$，所以
\begin{equation}
  |\arg(\mu)|=|\arg(-\lambda)|=\pi.
\end{equation}
又因为 $1<\beta<2$，有
\begin{equation}
  \pi>\frac{\beta\pi}{2}.
\end{equation}
因此 Matignon 稳定性条件成立。于是
\begin{equation}
  x_1(t)\to 0,\qquad
  x_2(t)=\dot{x}_1(t)\to 0.
\end{equation}
这说明到达滑模面后，分数阶滑动动力学本身是稳定的。

\subsection{Lyapunov 泛函证明}

上述 Matignon 判据给出了线性分数阶系统的代数稳定性结论。若希望用
Lyapunov 方法说明同一结论，需要注意：由于本例中
$1<\beta<2$，不能直接套用常见的 $0<q<1$ 情形下的引理
\[
  \frac{1}{2}{}_0^{\mathrm{C}}D_t^q x^2(t)
  \le
  x(t)\,{}_0^{\mathrm{C}}D_t^q x(t).
\]
更合适的做法是回到滑动模态的积分形式。令
\begin{equation}
  \mu=\beta-1=1-\alpha,\qquad 0<\mu<1.
\end{equation}
由 \eqref{eq:sliding-dynamics} 可写为
\begin{equation}
  \dot{x}_1(t)
  =
  -\lambda\,{}_0 I_t^\mu x_1(t).
  \label{eq:volterra-sliding}
\end{equation}
这与
${}_0^{\mathrm{C}}D_t^\beta x_1=-\lambda x_1$
具有相同的特征方程；若从 Caputo 方程直接出发，则需要把
$\dot{x}_1(0)$ 对应的初值项显式保留或通过滑动面起始时刻重置。

利用分数阶积分的扩散表示
\begin{equation}
  {}_0 I_t^\mu x_1(t)
  =
  \int_0^\infty \omega_\mu(\theta) z(\theta,t)\,d\theta,
  \label{eq:diffusive-integral}
\end{equation}
其中
\begin{equation}
  \omega_\mu(\theta)
  =
  \frac{\sin(\pi\mu)}{\pi}\theta^{-\mu}>0,
  \qquad
  \frac{\partial z(\theta,t)}{\partial t}
  =
  -\theta z(\theta,t)+x_1(t),
  \qquad z(\theta,0)=0.
\end{equation}
于是 \eqref{eq:volterra-sliding} 等价于扩展系统
\begin{equation}
  \dot{x}_1(t)
  =
  -\lambda
  \int_0^\infty \omega_\mu(\theta) z(\theta,t)\,d\theta.
  \label{eq:diffusive-system}
\end{equation}

选取 Lyapunov 泛函
\begin{equation}
  V(t)
  =
  \frac{1}{2}x_1^2(t)
  +
  \frac{\lambda}{2}
  \int_0^\infty
  \omega_\mu(\theta)z^2(\theta,t)\,d\theta.
  \label{eq:diffusive-lyapunov}
\end{equation}
由于 $\lambda>0$ 且 $\omega_\mu(\theta)>0$，有 $V(t)\ge 0$，并且
$V=0$ 当且仅当 $x_1=0$ 且 $z(\theta,t)=0$。对
\eqref{eq:diffusive-lyapunov} 沿 \eqref{eq:diffusive-system} 求导，得到
\begin{align}
  \dot V
  &=
  x_1\dot{x}_1
  +
  \lambda
  \int_0^\infty
  \omega_\mu(\theta)z(\theta,t)
  \frac{\partial z(\theta,t)}{\partial t}\,d\theta
  \notag\\
  &=
  -\lambda x_1
  \int_0^\infty \omega_\mu(\theta)z(\theta,t)\,d\theta
  +
  \lambda
  \int_0^\infty
  \omega_\mu(\theta)z(\theta,t)
  \bigl[-\theta z(\theta,t)+x_1\bigr]\,d\theta
  \notag\\
  &=
  -\lambda
  \int_0^\infty
  \theta\,\omega_\mu(\theta)z^2(\theta,t)\,d\theta
  \le 0.
  \label{eq:diffusive-vdot}
\end{align}
因此扩展状态是 Lyapunov 稳定的。进一步，由
\eqref{eq:diffusive-vdot} 可知 $\dot V=0$ 时有
$z(\theta,t)=0$ 几乎处处成立。在该不变集合内，
\[
  \frac{\partial z}{\partial t}
  =
  -\theta z+x_1=0
\]
推出 $x_1=0$。根据 LaSalle 不变性原理，得到
\begin{equation}
  x_1(t)\to 0.
\end{equation}
又由 \eqref{eq:volterra-sliding} 和 \eqref{eq:diffusive-integral}，
\begin{equation}
  x_2(t)=\dot{x}_1(t)
  =
  -\lambda\,{}_0 I_t^\mu x_1(t)
  \to 0.
\end{equation}
因此滑动阶段不仅满足 Matignon 判据，也可以通过上述
Lyapunov 泛函证明其渐近稳定性。

\section{采用饱和函数时的最终有界性}

若采用 \eqref{eq:control-sat}，当 $|s|>\phi$ 时，稳定性推导与符号函数情形一致：
\begin{equation}
  \dot V_s\le -k s^2-(\rho-\bar d)|s|.
\end{equation}
当 $|s|\le \phi$ 时，有
\begin{equation}
  \dot{s}
  =
  -\left(k+\frac{\rho}{\phi}\right)s+d(t).
\end{equation}
因此边界层内的滑模变量满足最终有界估计
\begin{equation}
  |s(t)|
  \le
  \frac{\bar d}{k+\rho/\phi}
  \quad \text{当 } t \text{ 足够大时}.
\end{equation}
可见，较小的 $\phi$、较大的 $k$ 或 $\rho$ 可以减小稳态边界，但过大的 $\rho$ 会增强控制输入的高频变化，因此工程实现中需要在鲁棒性和抖振之间折中。

\section{数值参数案例}

取一组参数
\begin{equation}
  \alpha=0.6,\qquad
  \lambda=2,\qquad
  k=3,\qquad
  \bar d=0.2,\qquad
  \rho=0.5.
\end{equation}
由于 $\rho=0.5>\bar d=0.2$，到达条件成立。分数阶滑模面为
\begin{equation}
  s=x_2+2\,{}_0 I_t^{0.4}x_1.
\end{equation}
理想符号控制律为
\begin{equation}
  u
  =
  -2\,{}_0^{\mathrm{RL}}D_t^{0.6}x_1
  -3s
  -0.5\,\sgn(s).
\end{equation}
若初始条件使得 $s(0)=1$，则有限时间到达上界为
\begin{equation}
  T_r
  \le
  \frac{1}{0.5-0.2}
  =
  3.33\ \mathrm{s}.
\end{equation}
滑动阶段的特征方程为
\begin{equation}
  p^{1.4}+2=0.
\end{equation}
等效标量分数阶系统为 ${}_0^{\mathrm{C}}D_t^{1.4}x_1=-2x_1$，其系统特征值为 $\mu=-2$。
由于
\begin{equation}
  |\arg(-2)|=\pi>0.7\pi=\frac{1.4\pi}{2},
\end{equation}
满足 Matignon 稳定性条件，因此滑动阶段稳定。

\section{结论}

本案例为双积分系统构造了分数阶滑模面
\[
  s=x_2+\lambda\,{}_0 I_t^{1-\alpha}x_1,
\]
并通过控制律抵消分数阶项，使滑模变量满足一阶鲁棒到达动态。Lyapunov 分析证明了当 $\rho>\bar d$ 时滑模面有限时间可达；进入滑模面后，系统满足阶次 $\beta=2-\alpha$ 的稳定分数阶动力学，因此状态收敛到原点。若用饱和函数替代符号函数，则系统不再严格到达 $s=0$，但可保证滑模变量最终有界。

\end{document}
